\documentclass[article]{jss}

\usepackage{amsmath}
\usepackage{amssymb}
\usepackage{graphicx}
\usepackage{url}
\usepackage{hyperref}
\usepackage{booktabs}
\usepackage{color}
\usepackage{natbib}

%% Fallbacks so the paper also builds against a minimal jss.cls (the real JSS
%% class defines all of these). Harmless when the full class is present.
\providecommand{\proglang}[1]{\textsf{#1}}
\providecommand{\pkg}[1]{\textbf{#1}}
\providecommand{\code}[1]{\texttt{#1}}

\setlength{\emergencystretch}{3em}

%% JSS article information
\author{
  Peter Cotton\\
  Microprediction\\
  \email{peter.cotton@microprediction.com}\\
  \url{https://humpday.microprediction.org}
}

\title{HumpDay: Derivative-Free Optimizers in Pure \proglang{Python} and
\proglang{JavaScript}, with a Contamination-Resistant Real-World Benchmark}

\Plaintitle{HumpDay: Derivative-Free Optimizers in Pure Python and
JavaScript, with a Contamination-Resistant Real-World Benchmark}
\Shorttitle{HumpDay}

\newcommand{\humpdayabstract}{
The \proglang{Python} package \pkg{humpday} provides twenty-three
derivative-free optimizers behind one uniform contract: minimise a black-box
function on the unit cube under a hard evaluation budget. Every algorithm is
implemented in pure \proglang{Python} with no required dependencies (a
\pkg{numpy} backend is used transparently when present) and again in
zero-dependency \proglang{JavaScript}, with parity tests holding the two in
agreement, so the same optimizer runs on a server or in a browser. A
recommender selects an algorithm from the problem's dimension, budget, and
measured evaluation cost, using rankings precomputed on the package's own
benchmark: eighty-one objectives ported from real applications, each paired
with an interactive browser demonstration. Scoring applies seeded smooth
bijections of the cube that relocate every optimum, so no method, human or
machine, can score by memorising solutions; the suite remains valid when the
candidates are written by language models. We describe the design and two
findings the benchmark has produced: optimizer rankings on synthetic test
functions correlate only moderately with rankings on the real-world suite
(Kendall $\tau$ between 0.42 and 0.49 across budgets), and the suite
supported the generation and out-of-sample validation of Alloy, a
machine-designed optimizer that now ships in the package.
}

\begin{document}

\maketitle

\begin{abstract}
\humpdayabstract
\end{abstract}

\noindent\textbf{Keywords:} derivative-free optimization, black-box
optimization, benchmarking, algorithm selection, benchmark contamination,
\proglang{Python}, \proglang{JavaScript}, \pkg{humpday}.

%% -------------------------------------------------------------------------
\section{Introduction} \label{sec:intro}

A practitioner with an expensive black-box objective and a budget of a few
hundred evaluations faces two problems. The first is having good optimizers
available without a compilation step or a dependency tree. The second is
knowing which one to use, because the answer changes with dimension, budget,
and the character of the objective.

Excellent algorithm collections exist: \pkg{scipy}~\citep{scipy2020},
\pkg{nevergrad}~\citep{nevergrad}, \pkg{NLopt}~\citep{nlopt},
\pkg{pymoo}~\citep{pymoo}, and the PRIMA re-implementations of Powell's
solvers~\citep{prima}. Benchmark platforms also exist, notably
COCO/BBOB~\citep{coco} and IOHprofiler~\citep{ioh}. \pkg{humpday} began as a
package that wrapped and ranked the collections; the version described here
inverts that design. Every algorithm is now a first-class pure-\proglang{Python}
port with no required dependencies, twinned in \proglang{JavaScript}, and the
rankings that drive its recommendations come from its own benchmark of
real-application objectives rather than from synthetic test functions.

Two properties distinguish the package from its neighbours. The first is the
dual implementation: parity tests hold the \proglang{Python} and
\proglang{JavaScript} implementations in agreement, so every algorithm, and
every benchmark problem, runs identically in a browser, which is where the
package's documentation lives. The second is contamination resistance.
Benchmark objectives derive from public material, and modern candidates may
be written by language models trained on that material; \pkg{humpday} scores
every run through a seeded smooth bijection of the cube that relocates the
optimum, so memorised solutions earn nothing. Section~\ref{sec:findings}
shows this is not a hypothetical concern: the suite has already been used to
generate, select, and validate an optimizer written entirely by a language
model.

Section~\ref{sec:usage} shows the package in use. Section~\ref{sec:design}
describes the design. Section~\ref{sec:benchmark} describes the benchmark
and its disguise mechanism. Section~\ref{sec:findings} reports the two
findings, and Section~\ref{sec:discussion} closes with scope and
limitations.

%% -------------------------------------------------------------------------
\section{Using the package} \label{sec:usage}

Installation is \code{pip install humpday}, with no required dependencies.
The everyday interface is a single call:

\begin{verbatim}
from humpday import minimize

def objective(x):
    return (x[0] - 0.3) ** 2 + abs(x[1] - 0.6)

result = minimize(objective, bounds=[(-1, 1), (0, 2)],
                  options={"n_trials": 200})
\end{verbatim}

Rectangular bounds are mapped to the unit cube internally. When
\code{method} is omitted, the algorithm is chosen by the recommender
(Section~\ref{sec:design}) from the dimension, the budget, and the measured
cost of one objective call. Passing \code{method="Alloy"} or any other
roster name overrides the choice.

The lower-level contract is deliberately plain. Every optimizer is a class
with the signature

\begin{verbatim}
optimizer = Alloy(objective, n_trials, n_dim)
best_value, best_x = optimizer.optimize()
\end{verbatim}

where \code{objective} takes a list of \code{n\_dim} floats in $[0,1]$ and
the budget \code{n\_trials} is a hard cap on objective calls. An ask/tell
interface drives any optimizer incrementally when the caller owns the
evaluation loop, and produces trajectories identical to the monolithic run.

The same roster is available in \proglang{JavaScript}:

\begin{verbatim}
const { Alloy } = require('humpday');
const opt = new Alloy(objective, 200, 2);
const { bestValue, bestX } = opt.optimize();
\end{verbatim}

%% -------------------------------------------------------------------------
\section{Design} \label{sec:design}

\subsection{One contract, two languages}

The twenty-three optimizers span six classical families (Powell-style trust
region, simplex and direct search, evolutionary and swarm methods,
annealing and other metaheuristics, model-based sampling, and pure
baselines) plus one machine-designed member
(Section~\ref{sec:findings}). All operate on $[0,1]^n$ under a hard
evaluation budget, and nothing else: no gradients, no Jacobians, no
constraint callbacks. Problems with other domains reach the cube through
transforms shipped with the package, including an invertible map between
the cube and the probability simplex, so portfolio and allocation problems
with weights summing to one are optimized by unmodified box-domain
algorithms.

The \proglang{Python} implementations read their array primitives from a
small shim that re-exports \pkg{numpy} when it is installed and a pure list
backend otherwise; algorithm code is identical either way. The
\proglang{JavaScript} implementations are ports of the same logic with no
dependencies at all. A parity suite runs both on identical problems and
fails if either implementation regresses, which keeps the browser
demonstrations evidential rather than decorative: what runs on the
documentation site is the algorithm being documented.

\subsection{The recommender}

Algorithm choice is data. Each optimizer carries a cost tier, and a
precomputed grid stores, for a lattice of (dimension, budget) cells, the
Borda mean rank of every algorithm across the benchmark suite. Given a
problem, \code{humpday.eligibility.recommend} filters by tier against the
measured cost of one objective call, then returns the eligible algorithm
with the best mean rank on the nearest cell. An opt-in cost-aware variant
penalises algorithms whose own overhead would dominate a cheap objective.
The recommendation is therefore reproducible arithmetic on published
benchmark results, not editorial opinion.

%% -------------------------------------------------------------------------
\section{The benchmark} \label{sec:benchmark}

The suite contains eighty-one objectives ported from real applications:
enzyme kinetics, wind-farm layout, battery dispatch, lens design, antenna
arrays, pool break shots, rocket landing, and others, spanning 2 to 100
dimensions. Each ships in both languages, and each has an interactive
browser page where every optimizer in the roster can be raced on the actual
problem.

Raw objectives of this kind have a weakness as benchmarks: they are built
from public material, so their solutions may be memorised, by people, by
tuned heuristics, and now by language models that write candidate code.
\pkg{humpday} therefore never scores a raw objective. A seeded smooth
bijection of the cube (a composition of coordinate warps) is applied inside
the objective, relocating the optimum while preserving the problem's
character; each seed yields a different disguised instance. A candidate
that knows where the answer used to live gains nothing, and a suite of
disguised instances measures search, not recall.

Scores are reported as regret normalised against a fixed panel of standard
optimizers run on the same instance: 0 means matching the best panel member
everywhere, 1 the worst. This makes results comparable across objectives
whose raw scales differ by orders of magnitude.

%% -------------------------------------------------------------------------
\section{Two findings} \label{sec:findings}

\subsection{Synthetic rankings transfer only moderately}

Do rankings earned on synthetic test functions predict rankings on real
problems? The package's full roster, plus a reference CMA-ES from
\pkg{nevergrad}, was ranked on a battery of synthetic functions of the
BBOB style, and independently on the disguised real-world suite, at three
budgets. Table~\ref{tab:tau} reports the rank correlations.

\begin{table}[t]
\centering
\begin{tabular}{lccc}
\toprule
budget & Kendall $\tau$ & $p$ & Spearman $\rho$ \\
\midrule
60  & 0.49 & 0.001 & 0.66 \\
120 & 0.49 & 0.001 & 0.67 \\
240 & 0.42 & 0.006 & 0.58 \\
\bottomrule
\end{tabular}
\caption{Rank correlation between the 22-optimizer leaderboard on synthetic
test functions and on the disguised real-world suite, per evaluation
budget.}
\label{tab:tau}
\end{table}

The correlation is positive and significant, and far from 1. Roughly, a
synthetic leaderboard explains under half the ordering that matters on the
real suite, and the disagreement grows with budget. Trust-region methods
that dominate smooth synthetic functions drop sharply on the real problems,
while robust direct-search methods rise. This is the empirical case for
recommending from real-application rankings, and a caution against
transplanting synthetic benchmark conclusions into practice.

\subsection{The suite supports algorithm discovery: Alloy}

Because the disguise mechanism makes scores meaningful even when candidates
are machine-written, the suite can drive search over algorithms themselves.
In a companion note~\citep{alloy2026}, a language model was prompted to
blend the mechanisms of five classical optimizers in stated proportions;
candidates were selected on disguised instances and the winner was validated
on twenty-nine objectives never used in any selection step. That winner,
Alloy, had the best mean rank at every budget from 60 to 480 evaluations
against six competitors including CMA-ES, winning 64\% to 77\% of 580
instances pairwise (sign-test $p < 10^{-10}$). Alloy now ships in the
package as an ordinary roster member, in both languages.

The companion note reports the failures alongside: most blends are
mediocre, single generations vary severalfold, and a tuned generic template
matched every in-sample score yet collapsed out of sample. The benchmark's
out-of-sample discipline, problems verified untouched against the history
of the suite, is what separated the finding from the noise.

%% -------------------------------------------------------------------------
\section{Discussion} \label{sec:discussion}

\pkg{humpday}'s scope is deliberately narrow: box-bounded, derivative-free,
serial, small-budget optimization. Users with gradients, constraints beyond
bounds and the simplex, or budgets of $10^5$ evaluations are better served
elsewhere, and the recommender will not pretend otherwise. Within scope, the
package offers a property the larger ecosystems do not: the algorithm you
read about, the algorithm you run in the browser demonstration, and the
algorithm you import are verifiably the same object, ranked on problems
that resist memorisation.

The benchmark's disguise mechanism addresses training-set contamination for
solution locations, not for problem structure: a model that has seen a
problem family will still recognise its shape. And the suite's problems,
however varied, are one curation; the rank-correlation finding cuts both
ways, cautioning against over-trusting any single suite, including this
one.

%% -------------------------------------------------------------------------
\section{Availability} \label{sec:availability}

\pkg{humpday} is MIT-licensed at
\url{https://github.com/microprediction/humpday}, installable with
\code{pip install humpday}, with the \proglang{JavaScript} roster published
as the \pkg{humpday} package on npm. Documentation, the interactive
demonstrations, and the working papers behind Section~\ref{sec:findings}
are at \url{https://humpday.microprediction.org}. Every number in this
paper is reproducible from run logs in the repository's \code{papers/}
directory.

\bibliographystyle{plainnat}
\bibliography{humpday}

\end{document}
